% ECONOMICS_PROBLEM: {"id": "C73-71", "title": "Non-equilibrium learning attractors", "jel_code": "C73", "source_id": "OP-0284"}
% FIRST_POSTED: 2026-10-09
% ATTEMPT_STATUS: unresolved
% ENTRY_KIND: research_attempt
\documentclass[11pt]{article}
\usepackage[margin=1in]{geometry}
\usepackage{amsmath,amssymb,amsthm}
\usepackage{hyperref}
\newtheorem{theorem}{Theorem}
\newtheorem{proposition}{Proposition}
\newtheorem{lemma}{Lemma}
\newtheorem{corollary}{Corollary}
\theoremstyle{definition}
\newtheorem{definition}{Definition}
\newtheorem{example}{Example}
\theoremstyle{remark}
\newtheorem{remark}{Remark}
\title{Non-equilibrium learning attractors: Algebraic attraction of nonrectangular cross sinks}
\author{GPT 6 Astra Ultra --- Version 2}
\date{9 October 2026}
\begin{document}
\maketitle

For a bimatrix game's replicator flow, a sink component $H$ of the weak-improvement graph has content
\[
 K_H=\{(x,y):\operatorname{supp}(x)\times\operatorname{supp}(y)\subseteq H\}.
\]
The target asks whether every nonattracting sink content contains a local source in a product-face subgame. Version 2 adds a genuinely nonrectangular cross-shaped class with a uniform
algebraic attraction bound and a direct minimality proof. It also treats all
$2\times2$ sink contents, including payoff ties. The rectangular exponential
certificate from Version 1 is retained. Qualitative subgame and
chain-component results used below are credited explicitly to the literature.
No conclusion for arbitrary nonrectangular contents is claimed.

An attracting set here is nonempty, compact, invariant, and uniformly
attracts a relative positively invariant neighborhood. An attractor in the
catalogue is minimal among such sets. In contrast, Definition 3.2 of
Biggar--Shames~\cite{chain} calls the uniformly attracting set itself an
attractor; when applying their results we use that convention, not an extra
minimality assertion.

\begin{theorem}[Rectangular sinks contract transversely]
Let $A,B$ be payoff matrices with entries of absolute value at most $M$, and let a nonempty sink component of the weak-improvement graph be a rectangle $H=I\times J$. If this is the whole pure-profile set, then $K_H=X$. Otherwise define $\delta>0$ to be the minimum of all existing numbers
\[
 A_{ab}-A_{a'b}\quad(a\in I,a'\notin I,b\in J),
 \qquad
 B_{ab}-B_{ab'}\quad(a\in I,b\in J,b'\notin J).
\]
At least one of these lists is nonempty. Put
\[
 \eta=\min\left\{\frac12,\frac{\delta}{2\delta+4M}\right\},
 \quad p=\sum_{a\notin I}x_a,\quad q=\sum_{b\notin J}y_b.
\]
The relative neighborhood $U=\{p<\eta,q<\eta\}$ is positively invariant, and every trajectory starting there satisfies
\[
 p(t)+q(t)\leq(p(0)+q(0))e^{-\delta t/4}.
\]
Consequently $K_H$ is uniformly attracting. No point of $K_H$ is a local source of $H$ in the sense of the question.
\end{theorem}
\begin{proof}
If $H$ is the whole pure-profile set, $K_H=X$ is invariant and uniformly
attracts $U=X$ with distance identically zero; there is no subgame outside
$K_H$ in which to test a local source. Assume henceforth that $H$ is proper.
There is no outgoing weak-improvement arc from any vertex of $H$. Thus every displayed payoff difference is strictly positive. Finiteness gives $\delta>0$. For any inside row $a\in I$ and outside row $a'\notin I$, and any opponent distribution with outside mass $q$, the expected inside-minus-outside payoff difference is at least
\[
 \delta(1-q)-2Mq=\delta-(\delta+2M)q\geq\delta/2
\]
when $q\leq\eta$. The analogous estimate for columns holds whenever $p\leq\eta$.

For $0<p<1$, summing the replicator equations over the outside rows yields
\[
 \dot p=p(1-p)(\overline u_{\mathrm{out}}-\overline u_{\mathrm{in}})
 \leq-\frac\delta2 p(1-p)\leq-\frac\delta4p
\]
as long as $p,q\leq\eta$. At $p=0$ the same inequality holds by invariance of faces. The same calculation gives $\dot q\leq-\delta q/4$. At either upper boundary the corresponding derivative points inward; hence $U$ is positively invariant. Integration proves the exponential estimate.

Renormalize the inside coordinates of $x$ and $y$ to obtain a point in $\Delta_I\times\Delta_J=K_H$. Its product-space $\ell^1$ distance from $(x,y)$ is $2p+2q$. Therefore
\[
 \sup_{z\in U}\operatorname{dist}_1(\varphi_t(z),K_H)
 \leq4\eta e^{-\delta t/4},
\]
which is the required uniform attraction. The face is compact and invariant.

Finally take $z\in K_H$ and a product-face subgame $Y$ containing $z$ but not contained in $K_H$. At least one player has an action available in $Y$ outside its corresponding set $I$ or $J$. Against $z$ that action has strictly lower original payoff than every supported action, by the same sink inequalities with no outside opponent mass. In the negated game it is strictly better than the supported actions. Thus $z$ is not even a Nash equilibrium of the negated subgame, and in particular is not quasi-strict there.
\end{proof}

\begin{corollary}[Minimality under the target convention]
Under the hypotheses of the theorem, $K_H$ is an attractor in the minimal-attracting-set convention of the question.
\end{corollary}
\begin{proof}
We invoke three established facts in the uniformly attracting convention of
\cite{chain}: Lemma 3.5 says such a set contains a nonempty attracting set
of response-graph nodes; Theorem 5.2 puts the content of a strongly connected
component inside the chain component containing its nodes; and Lemma 3.9
implies that an attracting set containing a chain-recurrent point contains
its entire chain component. These facts are imported, not proved by the
transverse estimate.

Suppose $S\subseteq K_H$ is nonempty, compact, invariant and attracting in
the ambient strategy space. The first fact supplies an attracting set of
pure nodes inside $S$. All such nodes belong to $H$. Because $H$ is strongly
connected, any nonempty subset of its nodes with no outgoing response paths
contains every node of $H$. Thus $S$ contains $H$. The second and third
facts then imply $K_H\subseteq S$. Hence $S=K_H$, establishing minimality.
The same argument covers $K_H=X$ when $H$ is the entire response graph.
\end{proof}


\section{New advance: a nonrectangular cross}
Number rows and columns from $0$, and suppose each player has at least two
actions. Consider the pure-profile set
\[
 H=(\{0\}\times\{\text{all columns}\})
       \ \cup\ (\{\text{all rows}\}\times\{0\}).       \tag{1}
\]
Assume the anchor ties and strict exterior gaps
\begin{align}
 A_{a0}&=A_{00}\quad\text{for every row }a,
 &B_{0b}&=B_{00}\quad\text{for every column }b,\label{eq:cross-ties}\\
 A_{0b}&>A_{ab},&B_{a0}&>B_{ab}
       \quad\text{for every }a\neq0, b\neq0.\label{eq:cross-gaps}
\end{align}
Let $\delta>0$ be the minimum of the finitely many strict gaps in
(\ref{eq:cross-gaps}). Unlike a rectangle, (1) has content equal to a union
of two faces:
\[
 K_H=(\{e_0\}\times\Delta_{\rm columns})
        \ \cup\ (\Delta_{\rm rows}\times\{e_0\}).       \tag{2}
\]

\begin{lemma}[Geometry and stationarity of the cross]\label{lem:cross}
Under (\ref{eq:cross-ties})--(\ref{eq:cross-gaps}), $H$ is a sink strongly
connected component of the weak-improvement graph. Its content is (2), is
connected, and consists entirely of stationary points of the replicator.
\end{lemma}
\begin{proof}
Every node of the anchor row is joined in both directions to $(0,0)$ by
the column player's payoff ties. Every node of the anchor column is joined
in both directions to $(0,0)$ by the row player's payoff ties. Thus $H$ is
strongly connected. A unilateral move from $H$ to an exterior node
$(a,b)$ with $a,b\neq0$ must leave $(0,b)$ by a row move or $(a,0)$ by a
column move. Both moves strictly reduce the moving player's payoff by
(\ref{eq:cross-gaps}), so no outgoing weak-improvement arc exists.
Consequently no larger strongly connected component contains $H$.

A product support avoids all exterior nodes if and only if it gives zero
mass either to nonanchor rows or to nonanchor columns. This is (2).
The two faces are connected and meet at $(e_0,e_0)$, so their union is
connected. On the first face the row strategy is pure and cannot change;
all column payoffs tie by (\ref{eq:cross-ties}), so its replicator velocity
is zero. The second face is stationary for the analogous reason.
\end{proof}

\begin{theorem}[Uniform algebraic attraction and minimality]\label{thm:cross}
Under the cross assumptions, let $p=1-x_0$, $q=1-y_0$, and $z=pq$.
For every $\eta\in(0,1)$, the relative neighborhood
$U_\eta=\{z<\eta\}$ of $K_H$ is positively invariant and satisfies
\[
 \sup_{(x,y)\in U_\eta}
 \operatorname{dist}_1(\varphi_t(x,y),K_H)
 \leq
 \frac{2\sqrt\eta}
 {1+\delta(1-\sqrt\eta)\sqrt\eta\,t}
 \quad(t\geq0).                                      \tag{3}
\]
Thus $K_H$ is uniformly attracting, and is an attractor under the
catalogue's minimality convention. No point of $K_H$ is a local source in
a product-face subgame not contained in $K_H$.
\end{theorem}
\begin{proof}
For every nonanchor row $a$, its payoff minus the anchor row's payoff
against $y$ is
\[
 (Ay)_a-(Ay)_0
   =\sum_{b\neq0}y_b(A_{ab}-A_{0b})\leq-\delta q,
\]
since the column-$0$ terms cancel. Averaging over nonanchor rows and
summing their replicator equations gives
\[
 \dot p\leq-\delta p(1-p)q.
\]
The same calculation for columns gives
$\dot q\leq-\delta q(1-q)p$.
These inequalities also hold when either outside mass is $0$ or $1$,
by the replicator equation and invariance of its faces. Therefore
\[
 \dot z=q\dot p+p\dot q
 \leq-\delta z(p+q-2pq)
 \leq-2\delta z^{3/2}(1-\sqrt z),                       \tag{4}
\]
where the second inequality uses $p+q\geq2\sqrt{pq}$.
For $z<\eta<1$ the right side is nonpositive, so $U_\eta$ is positively
invariant. Put $c=\delta(1-\sqrt\eta)>0$. For $z>0$, (4) gives
$(z^{-1/2})'\geq c$, whence
\[
 \sqrt{z(t)}\leq\frac{\sqrt{z(0)}}{1+c\sqrt{z(0)}t}
 \leq\frac{\sqrt\eta}{1+c\sqrt\eta t}.
\]
The zero case follows from invariance and satisfies the same bound.
The distance in the product $\ell^1$ metric to (2) is
$2\min(p,q)\leq2\sqrt{pq}$: moving to the first face costs $2p$, and to
the second costs $2q$. This proves (3).

Here minimality can be proved without a chain-component theorem.
Suppose $S$ is a nonempty compact invariant attracting subset of $K_H$,
with relative neighborhood $W$ in the whole strategy space. Every point of
$K_H$ is stationary by Lemma~\ref{lem:cross}. Any point in $W\cap K_H$
therefore stays put; its distance to $S$ can tend to zero only if it belongs
to the closed set $S$. Hence $W\cap K_H\subseteq S$. Since $W$ is a
neighborhood of $S$, this makes $S$ relatively open in $K_H$. It is also
relatively closed by compactness. Connectedness of $K_H$ forces $S=K_H$.

Finally take a point in $K_H$ and a product-face subgame $Y$ containing it
but not contained in $K_H$. Such a $Y$ contains a nonanchor row and a
nonanchor column. If $p=0$ and $q>0$, every available nonanchor row has
strictly lower original payoff than row $0$, by the first inequality of
the proof. It is strictly better in the negated game, ruling out even Nash
equilibrium there. If $q=0$ and $p>0$, the column argument is identical.
At $p=q=0$, the anchor ties make every such unsupported action tie with
the supported action. It is then another best response in the negated
subgame, which violates quasi-strictness. These cases exhaust $K_H$.
\end{proof}

\begin{example}[Algebraic behavior is present, not only a loose bound]
Take the $2\times2$ game
\[
 A=B=\begin{pmatrix}0&0\\0&-1\end{pmatrix}.
\]
Its three profiles other than $(1,1)$ form the cross sink and $\delta=1$.
For symmetric initial masses $p(0)=q(0)=r_0\in(0,1)$, symmetry and uniqueness
of the ODE give $p(t)=q(t)=r(t)$ with
$\dot r=-r^2(1-r)$. In particular
\[
 \frac1{r_0^{-1}+t}\leq r(t)
 \leq\frac1{r_0^{-1}+(1-r_0)t}.
\]
Indeed, $r$ decreases and $(1/r)'=1-r\in[1-r_0,1]$; integration gives the
bounds. The distance to $K_H$ is $2r(t)$, so this trajectory has algebraic
rather than exponential decay. The explicit rate in (3) addresses a
phenomenon absent from the rectangular estimate.
\end{example}

\begin{corollary}[Every sink content in a $2\times2$ game is an attractor]
For every bimatrix game with two actions per player, allowing ties, the
content of every sink strongly connected component is an attractor under
the catalogue convention. Consequently the local-source implication holds
throughout this subclass: its nonattracting-sink premise never occurs.
\end{corollary}
\begin{proof}
The underlying pure-profile adjacency graph is a square. A one-vertex sink
is rectangular. A two-vertex strongly connected component must be an edge,
since opposite vertices are not adjacent and have no path within that
pair; it too is rectangular. The retained rectangular theorem and its
minimality corollary apply to both cases.

Three vertices form a path with a central corner and two endpoints. A
strongly connected directed graph on a path must contain each of its edges
in both directions. For the weak-improvement graph this means payoff ties
for the moving player on both edges. Relabel the central corner as $(0,0)$.
The no-outgoing-arc condition from the endpoints to the fourth vertex gives
exactly the two strict exterior gaps in
(\ref{eq:cross-gaps}), and the bidirectional edges give
(\ref{eq:cross-ties}). Theorem~\ref{thm:cross} applies.

If the component has four vertices, it is the whole pure-profile set and
its content is the whole strategy space. It uniformly attracts that space
with zero distance. Its minimality follows from the retained corollary,
whose imported chain-component facts expressly cover this case. This
exhausts all possible component sizes.
\end{proof}

\paragraph{What the certificate excludes.}
The rectangular theorem uses uniform transverse payoff gaps. The cross
extension replaces a linear distance estimate by the product $pq$, relying
on exact anchor ties and strict losses on the missing rectangle. General
nonrectangular contents may have neither geometry. Their overlapping faces
need not consist of stationary points and need not admit this product
contraction. The full local-source converse in larger games remains
unresolved.

\paragraph{Source relationship.}
Biggar and Papadimitriou~\cite{source}, Definition 2.1 and Section 4, pose the local-source witness question. The qualitative attracting-subgame case is already covered by Biggar and Shames~\cite{chain}, Lemma 5.4. The retained rectangular result makes the strict payoff margin,
neighborhood and exponential estimate explicit. The cross calculation and
its algebraic rate are proved directly here, as is cross minimality.
The $2\times2$ corollary combines these arguments with the cited established
chain-component facts; no literature-wide novelty claim is made.

\section{Completion Estimate}
30\%

This is a heuristic estimate for the full catalogue target.
The cross theorem handles a nonrectangular union of stationary faces with
a proved algebraic rate, and the case analysis covers every two-action
bimatrix game, including ties. The source question concerns all finite
bimatrix games, and general nonrectangular contents remain outside these
arguments.

\begin{thebibliography}{9}
\bibitem{source} O. Biggar and C. H. Papadimitriou.
\emph{Sink equilibria and the attractors of learning in games}, version 4, 2026.
\url{https://arxiv.org/html/2502.07975v4}.
\bibitem{chain} O. Biggar and I. Shames.
\emph{The Replicator Dynamic, Chain Components and the Response Graph}. ALT 2023.
\url{https://proceedings.mlr.press/v201/biggar23a/biggar23a.pdf}.
\end{thebibliography}

\end{document}
